\documentclass[11pt]{article}
\usepackage[a4paper,margin=28mm]{geometry}
\usepackage{amsmath,amssymb,amsthm,mathtools}
\usepackage[T1]{fontenc}
\usepackage{lmodern}
\usepackage[colorlinks=true,linkcolor=blue,citecolor=blue,urlcolor=blue]{hyperref}
\newtheorem{theorem}{Theorem}
\newtheorem{lemma}{Lemma}
\newcommand{\T}{\mathbb T}
\newcommand{\Z}{\mathbb Z}
\newcommand{\R}{\mathbb R}
\newcommand{\norm}[1]{\left\lVert#1\right\rVert}
\newcommand{\ip}[2]{\left\langle#1,#2\right\rangle}
\DeclareMathOperator{\diver}{div}
\title{Failure of a mixed Sobolev estimate for scalar potentials}
\author{Kaiwen Jin\thanks{Contact: \href{mailto:yui.ui.siki@gmail.com}{yui.ui.siki@gmail.com}.}}
\date{}
\begin{document}
\maketitle

\begin{abstract}
We construct smooth real mean-zero functions on the two-dimensional unit
torus whose gradients stay bounded in $H^{-1/2}+L^1$, while their
$H^{1/2}+W^{1,1}$ norms tend to infinity. This disproves the mixed
potential estimate proposed by Bourgain and Brezis. The construction uses
a finite stationary reflected walk to produce a bounded vector field
with a uniformly bounded $H^{-1/2}$ divergence. Pairing this divergence
with a smoothed sheet potential gives the divergent lower bound.
\end{abstract}

\section{The estimate and the counterexample}

On $\T^d=\R^d/\Z^d$, with measure of total mass one, consider
\begin{equation}\label{eq:question}
 \norm{\psi}_{H^{1/2}+W^{1,1}}
 \leq C_d\norm{\nabla\psi}_{H^{-1/2}+L^1},
 \qquad \psi\in C^\infty(\T^d;\R),\quad \int_{\T^d}\psi=0.
\end{equation}
The constant is to depend only on the dimension, for every $d\geq2$.
This estimate appears in Bourgain--Brezis \cite[eq.~(1.16), p.~397]{BB}
and in Brezis \cite[Open Problem~7.1]{Brezis}.

Our Fourier convention and inhomogeneous Sobolev norms are
\[
 \widehat u(k)=\int_{\T^d}u(x)e^{-2\pi i k\cdot x}\,dx,
 \qquad
 \norm{u}_{H^s}^2=\sum_{k\in\Z^d}(1+4\pi^2|k|^2)^s|\widehat u(k)|^2.
\]
For vectors, the modulus is Euclidean. We use
$\norm v_{W^{1,1}}=\norm v_1+\norm{\nabla v}_1$, and
\[
 \norm w_{X+Y}=\inf_{w=u+v}
                  \bigl(\norm u_X+\norm v_Y\bigr),
\]
where equality in the decomposition is distributional. Using a sum of
component gradient norms in $W^{1,1}$ gives an equivalent norm with constants
depending only on $d$, and hence the same failure of \eqref{eq:question}.

\begin{theorem}\label{thm:main}
There are constants $c,C>0$ and real mean-zero
$\psi_M\in C^\infty(\T^2)$, for every integer $M\geq2$, such that
\begin{equation}\label{eq:main}
 \norm{\nabla\psi_M}_{H^{-1/2}(\T^2;\R^2)+L^1(\T^2;\R^2)}\leq C,
 \qquad
 \norm{\psi_M}_{H^{1/2}(\T^2)+W^{1,1}(\T^2)}\geq c\sqrt M-C.
\end{equation}
Consequently no finite $C_2$ satisfies \eqref{eq:question}.
\end{theorem}

The lower bound will come from functions $f_M=\diver Z_M$ with both
$\norm{f_M}_{H^{-1/2}}$ and $\norm{Z_M}_\infty$ bounded independently of $M$.
Such a function tests every decomposition of the scalar potential at once.
To obtain a large pairing, we first use a sheet supported at $x=0$,
then smooth it. The scales normal to the sheet are separated enough that
the divergence errors have bounded total Sobolev energy, whereas their
pairings add with the same sign.

We write the coordinates as $(x,y)$, with $x\in(-1/2,1/2)$ and
$y\in[0,1)$. Constants in the proof depend only on a cutoff fixed below.
They are independent of all integers and small scales in the construction.
We use the usual $L^2$ Fourier expansion on a torus
\cite[\S4]{Guillemin}; passing from period $2\pi$ to period one by
$\theta=2\pi x$ gives precisely the convention above. Every interchange
of integrals below involves absolutely integrable functions, so Fubini's
theorem applies \cite[Theorem~1.7.21]{Tao}.

\section{A finite reflected walk}

Fix $M\geq2$ and put
\begin{equation}\label{eq:scales}
 \beta=M^{-1},\qquad J=M^3,\qquad L=2M^6,
 \qquad h_0=(2M+1)^{-1},\qquad h_j=h_0L^{-j}.
\end{equation}
Partition the $y$-circle into $2M+1$ equal intervals, and assign the values
\[
 0,\ \beta,\ -\beta,\ 2\beta,\ -2\beta,\ \ldots,\ 1,\ -1
\]
in this order. Call this step function $U_0$. Refine every interval into
$L$ equal children, and repeat $J$ times. At refinement $j$, let $r_j$
be $+1$ on even children and $-1$ on odd children, starting with child zero.
Lebesgue measure makes $U_0$ uniform on the $2M+1$ states
$\{-1,-1+\beta,\ldots,1\}$ and the $r_j$ independent uniform signs,
independent of $U_0$. Indeed every parent has exactly $L/2$ children of
each sign. This is only a finite partition construction.

Define
\[
 U_j=\max\{-1,\min\{1,U_{j-1}+\beta r_j\}\}.
\]
The transition matrix is symmetric: at an interior state the walk moves
to either neighbor with probability $1/2$, while at a boundary state it
stays or moves inward, each with probability $1/2$. Therefore the uniform
law is stationary. Decompose the increments as
\begin{equation}\label{eq:increments}
 U_j-U_{j-1}=d_j+e_j,\qquad
 e_j=-\frac\beta2\operatorname{sgn}(U_{j-1})
                   \mathbf1_{\{|U_{j-1}|=1\}},\qquad d_j=c_jr_j,
\end{equation}
where $c_j=\beta$ in the interior and $c_j=\beta/2$ at the boundary.
Here $c_j,e_j$ are constant on each parent of length $h_{j-1}$, and
$d_j$ has mean zero on each parent. Set
\[
 a_M=J^{-1/2}\sum_{j=1}^Jr_j.
\]

\begin{lemma}\label{lem:walk}
The functions above satisfy $\int a_M=0$, $\norm{a_M}_2=1$, and
\begin{equation}\label{eq:walkbounds}
 \sum_{j=1}^J\norm{e_j}_2^2
   =\frac{J\beta^2}{2(2M+1)}\leq\frac14,
 \qquad
 -\sum_{j=1}^J\ip{a_M}{e_j}\geq\frac12\sqrt M-1.
\end{equation}
\end{lemma}
\begin{proof}
The first two assertions follow by independence. Stationarity gives
$|\{|U_{j-1}|=1\}|=2/(2M+1)$, proving the energy identity.
In $\ip{a_M}{d_j}$, every term $r_i$ with $i<j$ vanishes after averaging
$r_j$ conditional on the past. Terms with $i>j$ vanish by independence
of $r_i$ from $d_j$. Thus
\[
 \ip{a_M}{d_j}=J^{-1/2}\int c_j\geq\frac\beta{2\sqrt J},
 \qquad \ip{a_M}{U_0}=0.
\]
Since $U_J=U_0+\sum_j(d_j+e_j)$ and $|U_J|\leq1$,
\[
 -\sum_j\ip{a_M}{e_j}
   =\sum_j\ip{a_M}{d_j}-\ip{a_M}{U_J}
   \geq\frac\beta2\sqrt J-1=\frac12\sqrt M-1.
\]
\end{proof}

\section{A bounded field and its divergence}

Fix a smooth even cutoff $\eta$, equal to one on $[-1,1]$, zero for
$|x|\geq2$, with $0\leq\eta\leq1$ and $\eta'\leq0$ on the positive
half-line. For example, with $b(s)=e^{-1/s}$ for $s>0$ and zero otherwise,
take $\eta(x)=b(2-|x|)/(b(2-|x|)+b(|x|-1))$.
Set $\alpha_j(x)=h_j^{-1}\eta'(x/h_j)$ and regard these as smooth periodic
functions. Their supports $h_j<|x|<2h_j$ do not overlap, and
$2h_0\leq2/5<1/2$.

Let
\[
 p_0(y)=\int_0^yU_0(s)\,ds,\qquad p_j(y)=\int_0^yd_j(s)\,ds\quad(j\geq1).
\]
The order $0,\beta,-\beta,\ldots,1,-1$ for $U_0$ and the alternating
children give
\begin{equation}\label{eq:primitives}
 0\leq p_0\leq h_0,\qquad 0\leq p_j\leq\beta h_j.
\end{equation}
Each primitive is continuous and periodic: every positive and negative
pair cancels, and $d_j$ has zero integral on every parent.
Define
\begin{align*}
 Z_1(x,y)&=U_0(y)\eta(x/h_0)
         +\sum_{j=1}^J(U_j-U_{j-1})(y)\eta(x/h_j),\\
 Z_2(x,y)&=-\alpha_0(x)p_0(y)-\sum_{j=1}^J\alpha_j(x)p_j(y).
\end{align*}
On each transition annulus $Z_1$ is a convex combination of consecutive
values of $U_j$, and elsewhere it is one of these values or zero.
Thus $|Z_1|\leq1$. Nonoverlap and \eqref{eq:primitives} give
\begin{equation}\label{eq:Zbound}
 \norm Z_\infty\leq D_\eta:=\sqrt{1+\norm{\eta'}_\infty^2}.
\end{equation}
Only the $x$ derivative of $Z_1$ and the $y$ derivative of $Z_2$ enter
the divergence. The $U_0$ and $d_j$ terms cancel, leaving, in distributions,
\begin{equation}\label{eq:f}
 f_M:=\diver Z=\sum_{j=1}^J\alpha_j(x)e_j(y).
\end{equation}
For fixed $M$ this is a bounded $L^2$ function of mean zero.

\begin{lemma}\label{lem:negative}
There is $A_\eta<\infty$, independent of $M$, such that
$\norm{f_M}_{H^{-1/2}(\T^2)}\leq A_\eta$.
\end{lemma}
\begin{proof}
For $0<t\leq1/5$, write $\alpha_t=t^{-1}\eta'(x/t)$.
If $\widehat\eta_{\R}(u)=\int_\R\eta(v)e^{-iuv}\,dv$, then
\[
 \widehat\alpha_t(n)=i2\pi nt\,\widehat\eta_{\R}(2\pi nt),\qquad
 \widehat\alpha_t(0)=0.
\]
Three integrations by parts and the trivial low-frequency bound imply
\begin{equation}\label{eq:alpha}
 |\widehat\alpha_t(n)|\leq B_\eta
           \min\{|n|t,(|n|t)^{-2}\}\quad(n\ne0).
\end{equation}
For $0<t\leq s\leq1/5$, split positive $n$ at $1/s$ and $1/t$.
The summands $n^{-1}|\widehat\alpha_t(n)\widehat\alpha_s(n)|$ are bounded
in the three ranges by constant multiples of
$nts$, $t/(n^2s^2)$, and $1/(n^5t^2s^2)$, respectively. Summing gives
\begin{equation}\label{eq:overlap}
 \sum_{n\ne0}\frac{|\widehat\alpha_t(n)\widehat\alpha_s(n)|}{|n|}
       \leq C_\eta\frac ts.
\end{equation}
The last range costs $C_\eta(t/s)^2\leq C_\eta t/s$; rounding the two
cutoffs changes only a constant.

For arbitrary $v_j\in L^2(\T)$, Parseval in $y$ and
$(1+4\pi^2(n^2+m^2))^{-1/2}\leq(2\pi|n|)^{-1}$ for $n\ne0$ give
\begin{align*}
 \norm{\sum_j\alpha_jv_j}_{H^{-1/2}}^2
 &\leq\frac1{2\pi}\sum_{n\ne0}\frac1{|n|}
            \norm{\sum_j\widehat\alpha_j(n)v_j}_2^2\\
 &\leq C_\eta\sum_{j,\ell}L^{-|j-\ell|}\norm{v_j}_2\norm{v_\ell}_2
 \leq C'_\eta\sum_j\norm{v_j}_2^2.
\end{align*}
The final bound uses $2ab\leq a^2+b^2$ and the row sum
$1+2/(L-1)\leq3$. Apply this to $v_j=e_j$ and use
Lemma~\ref{lem:walk}. No tangential orthogonality is required.
\end{proof}

\section{The sheet potential}

Let $P_t$ be the circle Poisson operator, with Fourier multiplier
$e^{-2\pi t|m|}$. The geometric series gives its nonnegative unit-mass kernel
\begin{equation}\label{eq:poisson}
 q_t(y)=\frac{1-r^2}{1-2r\cos(2\pi y)+r^2},\qquad r=e^{-2\pi t}.
\end{equation}
For $0<t\leq1$ it satisfies
$q_t(y)\leq Ct/(t^2+\operatorname{dist}(y,\Z)^2)$.

\begin{lemma}\label{lem:poisson}
If $e$ is constant on equal circle intervals of length $h$ and takes
values in $[-\beta/2,\beta/2]$, then
\[
 \norm{P_te-e}_2\leq C\beta\sqrt{t/h},\qquad 0<t\leq h.
\]
In particular, for $h_j\leq x\leq2h_j$,
$\norm{P_xe_j-e_j}_2\leq C\beta/\sqrt L$.
\end{lemma}
\begin{proof}
Translation by distance $s\leq1/2$ changes $e$ only on a set of measure
at most $\min(1,s/h)$, and the change is at most $\beta$. Thus the
$L^2$ difference is at most $\beta\min(1,\sqrt{s/h})$.
The integral triangle inequality and the kernel bound reduce the claim
to integrating $Ct/(t^2+s^2)\min(1,\sqrt{s/h})$.
For $0<s<h$, substitution $s=tv$ bounds the integral by
\[
 C\sqrt{t/h}\int_0^\infty\frac{\sqrt v}{1+v^2}\,dv.
\]
For $s>h$ the bound is $Ct/h\leq C\sqrt{t/h}$.
The two directions on the circle change only the constant.
Finally $e_j$ is constant on parents of length $h_{j-1}$, and
$x\leq2h_j=2h_{j-1}/L$.
\end{proof}

Temporarily take the vector distribution
$F_0=(1,0)\delta_0(x)a_M(y)$, and define the mean-zero potential
$\psi_0=\Delta^{-1}\diver F_0$. Its nonzero Fourier coefficients are
\begin{equation}\label{eq:sheetfourier}
 \widehat\psi_0(n,m)=-\frac{in\widehat a_M(m)}{2\pi(n^2+m^2)}.
\end{equation}
All $m=0$ coefficients vanish. Since
\[
 \sum_n\frac{n^2}{(n^2+m^2)^2}
       \leq\sum_n\frac1{n^2+m^2}\leq\frac C{|m|}\quad(m\ne0),
\]
we have $\psi_0\in L^2(\T^2)$ with norm bounded independently of $M$.

For $0<x<1/2$ and $\lambda=2\pi|m|$, the periodic Green function for
$\partial_x^2-\lambda^2$ is
\[
 G(x)=-\frac{e^{-\lambda x}+e^{-\lambda(1-x)}}
                     {2\lambda(1-e^{-\lambda})},\qquad 0<x<1.
\]
It is continuous across the identified endpoints and $G'$ has jump one.
Differentiating therefore gives
\begin{equation}\label{eq:green}
 \widehat{\psi_0(x,\cdot)}(m)
       =\frac{e^{-\lambda x}-e^{-\lambda(1-x)}}
                   {2(1-e^{-\lambda})}\widehat a_M(m).
\end{equation}
The potential is odd in $x$. On $0<x\leq2/5$ write this as
$\psi_0(x,\cdot)=\tfrac12P_xa_M+R_xa_M$.
On mean-zero $L^2(\T)$ functions, the multiplier of $R_x$ is
$-e^{-\lambda}\sinh(\lambda x)/(1-e^{-\lambda})$.
Using $\sinh(\lambda x)\leq\lambda x e^{\lambda x}$ and
$\lambda\geq2\pi$ shows $\norm{R_x}_{L^2\to L^2}\leq Cx$.
Self-adjointness of $P_x$, Lemma~\ref{lem:poisson}, and $\norm{a_M}_2=1$
now give, on the positive support of $\alpha_j$,
\[
 \left|\ip{\psi_0(x,\cdot)}{e_j}-\frac12\ip{a_M}{e_j}\right|
       \leq C\beta/\sqrt L+Cx\norm{e_j}_2.
\]
Both $\psi_0$ and $\alpha_j$ are odd in $x$, and
$\alpha_j\leq0$ for $x>0$ with $\int_{x>0}\alpha_j=-1$.
Integrating, summing, and using \eqref{eq:scales} yields
\begin{equation}\label{eq:pairing}
 \left|\int_{\T^2}\psi_0f_M+\sum_j\ip{a_M}{e_j}\right|
 \leq CJ\beta/\sqrt L+C\sum_jh_j\norm{e_j}_2\leq C/M.
\end{equation}
Indeed $J\beta/\sqrt L=1/(\sqrt2M)$ and
$\sum_jh_j\norm{e_j}_2\leq\beta h_1/(1-1/L)$.
All these are legitimate $L^2$ pairings for each finite $M$.
Lemma~\ref{lem:walk} proves
\begin{equation}\label{eq:largepair}
 \int\psi_0f_M\geq\frac12\sqrt M-1-C/M.
\end{equation}

\section{Smoothing and the mixed norms}

For $\delta>0$ let $\rho_\delta=q_\delta$ from \eqref{eq:poisson}, and set
\[
 a_{M,\delta}=\rho_\delta*a_M,\qquad
 F_{M,\delta}(x,y)=(1,0)\rho_\delta(x)a_{M,\delta}(y),\qquad
 \psi_{M,\delta}=\Delta^{-1}\diver F_{M,\delta}.
\]
Their Fourier series have exponential decay. Thus the functions are smooth
and real, and the potential has mean zero. Positivity and unit mass give
$\norm{F_{M,\delta}}_1\leq\norm{a_M}_1\leq1$.

For $k=(n,m)\ne0$ put $P(k)=kk^T/|k|^2$, and put $P(0)=0$, $Q=I-P$.
The Fourier formula for the inverse Laplacian gives
$\nabla\psi_{M,\delta}=F_{M,\delta}-QF_{M,\delta}$.
Since $|Q(n,m)(1,0)|^2=m^2/(n^2+m^2)$ and all $m=0$ terms vanish,
\begin{align*}
 \norm{QF_{M,\delta}}_{H^{-1/2}}^2
 &\leq\sum_{m\ne0}|\widehat a_{M,\delta}(m)|^2
     \sum_n\frac{m^2}{(n^2+m^2)\sqrt{1+4\pi^2(n^2+m^2)}}\\
 &\leq\frac3{2\pi}\norm{a_{M,\delta}}_2^2\leq\frac3{2\pi}.
\end{align*}
For the middle bound, the decreasing even function
$m^2/(t^2+m^2)^{3/2}$ has integral one on $t>0$, so its integer sum
is at most $1/|m|+2\leq3$. We have proved the uniform estimate
\begin{equation}\label{eq:right}
 \norm{\nabla\psi_{M,\delta}}_{H^{-1/2}+L^1}
       \leq B:=1+\sqrt{3/(2\pi)}.
\end{equation}

For fixed $M$, smoothing multiplies \eqref{eq:sheetfourier} by
$e^{-2\pi\delta(|n|+|m|)}$, so $\psi_{M,\delta}\to\psi_0$ in $L^2$.
To see this without any uniformity assertion in $M$, first bound the
squared error outside a finite set by four times the Fourier tail of
$\psi_0$, then let $\delta$ decrease on the finite set.
Let $n_M$ be the least integer $n\geq M$ for which
\[
 \norm{\psi_{M,2^{-n}}-\psi_0}_2\leq(1+\norm{f_M}_2)^{-1}.
\]
It exists by this convergence. Set $\delta_M=2^{-n_M}$ and
$\psi_M=\psi_{M,\delta_M}$. Then \eqref{eq:largepair} implies
\begin{equation}\label{eq:smoothpair}
 \int\psi_Mf_M\geq\frac12\sqrt M-2-C/M.
\end{equation}

For any decomposition $\psi=u+v$ with $u\in H^{1/2}$ and
$v\in W^{1,1}$, Fourier Cauchy--Schwarz and \eqref{eq:f} give
\begin{align*}
 \left|\int\psi f_M\right|
 &\leq\norm{f_M}_{H^{-1/2}}\norm u_{H^{1/2}}
              +\norm Z_\infty\norm{\nabla v}_1\\
 &\leq C_*\bigl(\norm u_{H^{1/2}}+\norm v_{W^{1,1}}\bigr),
 \qquad C_*:=\max(A_\eta,D_\eta).
\end{align*}
For the integration by parts, extend $v$ periodically and mollify on an
open cube larger than a fundamental cube. Local weak-derivative
mollification \cite[Theorem~3.19]{Hunter} gives smooth periodic approximations
converging with their gradients in $L^1(\T^2)$.
The integrals against fixed-$M$ bounded $f_M$ and bounded $Z$ therefore
converge. Also $u\in L^2\subset L^1$, so a distributional decomposition
agrees with an almost-everywhere equality and the displayed pairing is valid.
Taking the infimum over all decompositions and using \eqref{eq:smoothpair},
\[
 \norm{\psi_M}_{H^{1/2}+W^{1,1}}
       \geq\frac{\frac12\sqrt M-2-C/M}{C_*}.
\]
Together with \eqref{eq:right}, this proves Theorem~\ref{thm:main} and
shows that the proposed family of estimates fails already in dimension two.

\paragraph{AI use.}
OpenAI Codex was used to develop the construction, draft the proof, check
sources, and prepare this manuscript. Independent mathematical review is
requested.

\begin{thebibliography}{9}
\bibitem{BB}
J. Bourgain and H. Brezis,
\emph{On the equation $\diver Y=f$ and application to control of phases},
Journal of the American Mathematical Society \textbf{16} (2003), 393--426.
\href{https://sites.math.rutgers.edu/~brezis/PUBlications/171-J.pdf}{Author PDF}.
\bibitem{Brezis}
H. Brezis, \emph{Some of my favorite open problems},
Rendiconti Lincei -- Matematica e Applicazioni \textbf{34} (2023),
no. 2, 307--335.
\href{https://doi.org/10.4171/RLM/1008}{doi:10.4171/RLM/1008}.
\bibitem{Guillemin}
V. Guillemin, \emph{Notes for 18.117: Elliptic operators}, MIT, Spring 2005,
\S4. \href{https://math.mit.edu/~vwg/classnotes-spring05.pdf}{Author notes}.
\bibitem{Tao}
T. Tao, \emph{An Introduction to Measure Theory}, Graduate Studies in
Mathematics 126, American Mathematical Society, 2011.
\href{https://terrytao.wordpress.com/wp-content/uploads/2012/12/gsm-126-tao5-measure-book.pdf}{Author's preliminary version}.
\bibitem{Hunter}
J. K. Hunter, \emph{Notes on Partial Differential Equations}, Chapter 3,
Theorem 3.19. \href{https://www.math.ucdavis.edu/~hunter/pdes/ch3.pdf}{Author notes}.
\end{thebibliography}
\end{document}
